% !TEX program = xelatex
% 공통수학2 · Ⅲ. 함수와 그래프 · 1. 함수 · 중단원 마무리 (38차시)
\documentclass[11pt,a4paper]{article}
\usepackage{kotex}
\usepackage{iftex}
\ifPDFTeX\else
  \setmainhangulfont{Noto Serif CJK KR}
  \setsanshangulfont{Noto Sans CJK KR}
\fi
\usepackage[margin=2.1cm]{geometry}
\usepackage{parskip}
\usepackage{enumitem}
\usepackage{array}
\usepackage{tabularx}
\usepackage{booktabs}
\usepackage{xcolor}
\usepackage{amssymb}
\usepackage{tikz}
\usepackage[most]{tcolorbox}
\usepackage{fancyhdr}
\usepackage{titlesec}

\definecolor{goldc}{HTML}{B07C22}
\definecolor{goldstrong}{HTML}{8A5F16}
\definecolor{indigoc}{HTML}{2B3B57}
\definecolor{indigosoft}{HTML}{6E82A6}
\definecolor{linec}{HTML}{D6DACB}
\definecolor{surface2}{HTML}{E4E8DC}
\definecolor{notebg}{HTML}{FBF3E3}
\definecolor{inksoft}{HTML}{52604F}

\titleformat{\section}{\normalfont\sffamily\bfseries\large\color{indigoc}}{\thesection}{0.6em}{}
\titlespacing{\section}{0pt}{16pt}{8pt}
\newtcolorbox{ideabox}{enhanced, breakable, colback=white, colframe=goldc, boxrule=0.5pt, leftrule=3pt, arc=2pt, left=10pt, right=10pt, top=8pt, bottom=8pt}
\newtcolorbox{calloutbox}{enhanced, breakable, colback=white, colframe=indigosoft, boxrule=0.5pt, leftrule=3pt, arc=2pt, left=10pt, right=10pt, top=7pt, bottom=7pt}
\newtcolorbox{answerbox}{enhanced, breakable, colback=white, colframe=linec, boxrule=0.5pt, arc=2pt, left=10pt, right=10pt, top=7pt, bottom=7pt}
\newcommand{\lessonbanner}[3]{%
  \begin{tcolorbox}[enhanced, colback=indigoc, colframe=indigoc, boxrule=0pt, arc=0pt,
    left=14pt, right=14pt, top=14pt, bottom=10pt, borderline south={2.4pt}{0pt}{goldc}, fontupper=\color{white}]
    {\sffamily\footnotesize\color{white!85!black} #1}\\[4pt]{\sffamily\bfseries\Large #2}\\[3pt]{\sffamily\small\color{white!88!black} #3}
  \end{tcolorbox}}
\pagestyle{fancy}\fancyhf{}\renewcommand{\headrulewidth}{0pt}
\fancyfoot[L]{\footnotesize\color{inksoft} 공통수학2 · 2026학년도 2학기 · 지평선고등학교}
\fancyfoot[R]{\footnotesize\color{inksoft} 교사용 지도안 -- 배포 금지}

\begin{document}
\lessonbanner{공통수학2 · Ⅲ. 함수와 그래프}{1. 함수 --- 중단원 마무리 (38차시)}{교사용 지도안 (2026학년도 2학기)}
\vspace{10pt}

\noindent\begin{tabularx}{\linewidth}{@{}>{\bfseries\color{indigoc}}p{2.6cm}X@{}}
\toprule
대상 & 고1 · 2개 반 · 반당 13명 \\
차시 & 50분 (도입 5 · 전개 35 · 정리 10) \\
목적 & 33$\sim$37차시(함수의 뜻, 여러 가지 함수, 합성함수, 역함수)를 종합 정리하고 형성평가로 점검한다. \\
\bottomrule
\end{tabularx}

\section{개념 지도 총정리}
\begin{ideabox}
33$\sim$37차시 전체를 관통하는 흐름: \textbf{대응 중 존재성·유일성을 만족하는 것이 함수 $\to$ 대응의 성질(일대일성, 전사성)에 따라 함수를 분류 $\to$ 두 함수를 이어 붙이면 합성함수 $\to$ 일대일대응인 함수만 대응을 거꾸로 뒤집어 역함수를 만들 수 있다}.
\end{ideabox}
\begin{tabularx}{\linewidth}{@{}>{\bfseries\small}p{5cm}X@{}}
\toprule
\textbf{개념} & \textbf{핵심 내용} \\
\midrule
함수의 뜻 & $X$의 모든 원소에 $Y$의 원소가 빠짐없이(존재성), 오직 하나씩(유일성) 대응 \\
정의역·공역·치역 & $X$: 정의역, $Y$: 공역, $f(X)=\{f(x)\mid x\in X\}$: 치역 (치역 $\subset$ 공역) \\
일대일함수 / 일대일대응 & $x_1\neq x_2\Rightarrow f(x_1)\neq f(x_2)$ / 일대일함수이면서 치역$=$공역 \\
항등함수 / 상수함수 & $f(x)=x$ / $f(x)=c$ (상수) \\
합성함수 & $(g\circ f)(x)=g(f(x))$, 일반적으로 $g\circ f\neq f\circ g$, 결합법칙 성립 \\
역함수 & $f$가 일대일대응일 때 존재. $y=f(x)\iff x=f^{-1}(y)$, $(g\circ f)^{-1}=f^{-1}\circ g^{-1}$ \\
역함수의 그래프 & $y=f(x)$와 $y=f^{-1}(x)$는 직선 $y=x$에 대하여 대칭 \\
\bottomrule
\end{tabularx}

\section{수업 흐름}
\begin{tabularx}{\linewidth}{@{}>{\centering\arraybackslash}p{1.6cm}X>{\raggedright\arraybackslash\small}p{3.1cm}@{}}
\toprule
\textbf{단계} & \textbf{활동 \& 발문} & \textbf{ATL / VTR} \\
\midrule
\textcolor{goldstrong}{\textbf{도입}}\newline\footnotesize 5분 &
\textbf{마음공부 (2분)} --- 중단원 전체 유무념 대조.\newline
\textbf{개념 지도 확인 (3분)} --- 위 표 빈칸형으로 복습. &
ATL: 자기관리(복습 전략) \\
\addlinespace
\textcolor{goldstrong}{\textbf{전개}}\newline\footnotesize 35분 &
\textbf{기본 (15분)} --- 마무리 문제 01$\sim$04(정의역·공역·치역, 일대일함수 판별, 합성함수 계산, 역함수 구하기) 개별 풀이 후 짝 교차 채점.\newline
\textbf{표준 (10분)} --- 05, 06(역함수의 그래프 교점, $(g\circ f)^{-1}=f^{-1}\circ g^{-1}$ 검증)을 모둠 협력으로 해결.\newline
\textbf{도전 (10분)} --- 서술형 07, 08과 발전 09(항등함수·상수함수의 순서쌍 표현, 자기 자신이 역함수인 조건, 두 가지 방법 비교) 중 학급 수준에 맞게 선택하여 시범 풀이. &
ATT: 촉진자·평가자\newline ATL: 협업·자기점검 \\
\addlinespace
\textcolor{goldstrong}{\textbf{정리}}\newline\footnotesize 10분 &
\textbf{자기 평가 (5분)} --- 성취수준 체크리스트.\newline
\textbf{성찰 (5분)} --- 감사 일기 + 다음 중단원(유리함수와 무리함수) 예고. &
마음공부 · 자기평가 \\
\bottomrule
\end{tabularx}

\begin{calloutbox}
\textbf{지도상의 유의점} --- 오답이 나온 문제가 33$\sim$37차시 중 어느 차시 개념인지 학생 스스로 표시하게 하여 개인별 보충 지점을 시각화한다. 특히 일대일함수와 일대일대응, $g\circ f$와 $f\circ g$, $(g\circ f)^{-1}$와 $f^{-1}\circ g^{-1}$의 순서처럼 `헷갈리기 쉬운 짝'을 짝지어 다시 짚어 준다.
\end{calloutbox}

\section{형성평가 답안 (지도서 원문 01$\sim$09번)}
\begin{tabularx}{\linewidth}{@{}>{\bfseries\color{goldstrong}}p{1.4cm}X@{}}
\toprule
01 & $X=\{1,2,3,4\}$, $Y=\{1,\ldots,10\}$, $f(x)=2x$: 정의역 $=\{1,2,3,4\}$, 공역 $=\{1,\ldots,10\}$, 치역 $=\{2,4,6,8\}$ \\
02 & 일대일함수: ① $f(x)=5-x$, ④ $f(x)=x^3$ (②, ③은 아님) \\
03 & $f(x)=3x+2$, $g(x)=-x+1$: $(g\circ f)(1)=-4$, \; $(f\circ g)(1)=2$ \\
04 & $f(x)=4x-3$의 역함수: $f^{-1}(x)=\dfrac{x+3}{4}$ \\
05 & $f(x)=2x-6$과 역함수 그래프의 교점: $(6,6)$ (증가함수이므로 $f(x)=x$로 풀이) \\
06 & $f(x)=x+3$, $g(x)=2x$: $(g\circ f)^{-1}(x)=\dfrac{x-6}{2}=f^{-1}(g^{-1}(x))$로 성질 검증 \\
07 & $X=\{1,2,3\}$: 항등함수 $\{(1,1),(2,2),(3,3)\}$, 상수함수(값 3) $\{(1,3),(2,3),(3,3)\}$ \\
08 & $f(x)=ax+b$가 $f=f^{-1}$가 되는 조건: $a=-1$($b$는 임의의 실수) 또는 자명한 경우 $a=1,\,b=0$ \\
09 & $f(x)=2x-1$, $g(x)=x+3$: $(f\circ g)^{-1}(x)=\dfrac{x-5}{2}$ (직접 구한 값과 $g^{-1}\circ f^{-1}$로 구한 값이 일치) \\
\bottomrule
\end{tabularx}

\section{다음 중단원}
\begin{calloutbox}
\textbf{2. 유리함수와 무리함수 (39차시$\sim$).} 분수식·무리식으로 표현되는 함수의 정의역과 그래프를 다룬다. 지금까지 배운 함수·역함수의 개념이 그대로 이어진다.
\end{calloutbox}
\end{document}
